% Nouveau EGA V (cote 25): archive cote 25, site page 53 (PDF page 54).
% typescript. Draft compared: none.
% AI-transcribed by Claude Opus 5.5 (claude-opus-5-5, adaptive thinking, effort high), Message Batches API, batch ega25b (msgbatch_01A7n5YuXubUPNK1zsepi8RG), 2026-11-13; not yet checked by eye.
\scanpage{25/53}
\margin{\unsure{46}}

$\overset{\hand{\vee}}{P}$ (correspondant à la question d'étudier les sections hypérplanes de X
passant par un point donné x, ou tangent à X en un point \deleted{simple}\inserted{lisse} donné),
form\over{e}{é} donc des hyperplans contenant une sous-variété linéaire $L^{o}$ de P
(resp. un point,  ou l'image d'un espace tangent de X en un point lisse,
dans les deux cas envisagés), et on se pose la question de savoir si
pour le point générique \hand{$\zeta$} de L  (donc pour tous les points d'un ouvert
non vide de L),  $Y_{\hand{\zeta}}$  est \xout{\unsure{géomxxintègrexxxsurxdeuxmoins}} géom irréd., de
\margin{/} dim  dim X -1.   C'est là une variante  de Bertini, qui devrait figurer
dans  le n°3  et se traite par exactement la même méthode\inserted[typed]{ou, si on veut, s'y ramène.}  La question
dimensionnelle s'exprime simplement par  $f(X)\not\subset L^{o}$ , \xout{\unsure{Ilxyxaxdonexxxxx}}
i.e.  $X' =$ \xout{Xxxx} $f^{-1}(P - L^{o})$   est un ouvert dense de X, \xout{\unsure{etxlaxquestion}}
Soit Q l'espace projectif des hypérplans passant par $L^{o}$  (NB si  $L^{o}$ est
défini par un sous-espace vectoriel $F^{o}$ de E,  on a   $Q = P($\xout{E/}$F^{o})$ ) et
considérons le morphisme  canonique  (déduit de \xout{ExxxE/$F^{o}$} $F^{o}\longrightarrow E$ , cf
Chap II)
\[ u :\ P - L^{o} \longrightarrow Q, \]
et considérons
\[ g = u f' :\quad f^{-1}(P - L^{o}) = X' \longrightarrow Q , \]
alors \xout{\unsure{laxfamillexdesxxxX'xxxxxtesxxautr}}  $L \simeq \overset{\hand{\vee}}{Q}$ ,  et la famille des
$X'_{\hand{\zeta}}$  $(\hand{\zeta}\hand{\in} L)$  n'est autre que la famille des ``sections hyperplanes'' rela-%
tif à ce morphisme  g . D'autre part, on voit aussitôt que pour tout
$\hand{\zeta}\in L$  ``général''  $X'_{\hand{\zeta}}$ est dense dans $X_{\hand{\zeta}}$ , donc  $X'_{\hand{\zeta}}$  est géom irréd sss
X  l'est.  Ceci posé,  le théorème de Bertini-Zariski nous montre que
nous avons la conclusion voulue \xout{dix} d'irréductibilité pourvu que
dim $g(X') \geqslant 2$. (A vrai dire, tant qu'à faire on pourrait signaler une
réciproque dans 3.1. , savoir : \xout{Yxxestxgé} si X est géom irréd, Y  est
géom irréd sss  ou bien  dim $f(X)$; $\neq$ \unsure{2|1}, ou bien  dim $f(X)=1$ , et $f(X)$
est contenu dans une droite\hand{D} (définie sur $\bar{k}$ )\inserted{et la fibre générique de $X\to D$ est géom. irréd}\hand{)}; cela permet également

\mspage{53}
